\section{Discussion: The Results in Light of Recent Work}\label{sec:discussion}

Read against the literature of Section~\ref{sec:literature}, the results support three conclusions and point to
three extensions.

\paragraph{The counter-theoretical signs are the expected symptom of an incomplete information set.} The
overlap-region estimates imply that rate increases are followed by stronger activity. That is what one would see if
the Committee tightened when it had reason to expect strength that the econometrician's covariates do not capture.
This is the mechanism behind the information effect of \citet{nakamura2018} and, in the account of
\citet{bauer2023aer}, behind the predictability of market surprises by economic news. The finding that
Greenbook forecasts remove the price puzzle but not the real-activity puzzle fits the evidence in \citet{aruoba2025}
that the staff's written analysis contains information not summarised by its numerical forecasts. It also fits the
argument of \citet{romer2023} that isolating policy changes unrelated to the outlook requires reading the record,
not only conditioning on forecasts. The menus do not solve this problem, and they are not designed to. What they do
is separate it cleanly from the overlap problem, so that the remaining bias can be attributed to unconfoundedness.

\paragraph{The menu restriction is a genuine source of identifying information.} Imposing the menus raises the
likelihood of the decision model by as much as market expectations do, and it is exact rather than approximate.
Recent work that uses the same documents for other purposes \citep{doh2026,laarits2025} confirms that the Bluebook
alternatives are a well-defined and informative record of the Committee's deliberations. In the terms of the overlap
literature, the menus turn a problem of extreme estimated scores \citep{khan2010,heiler2021} into one of
structural zeros in a known region. That region can be handled by partial identification
\citep{susmann2025} instead of trimming, which would change the estimand.

\paragraph{Worst-case bounds alone are too wide for macroeconomic questions.} With one third to two fifths of
meetings outside the overlap region and outcome ranges that include the 2008--09 recession, the identified sets
contain zero everywhere. The same arithmetic makes the cross-sectional bounds of \citet{susmann2025} informative
only when the non-overlap subpopulation is small. In monetary policy it is not, so informative conclusions require
an additional assumption. The breakdown values $\lambda^\ast$ make that assumption explicit and give it units.

\paragraph{Extensions.} First, the menu-constrained propensity score can be combined with a richer information set.
Text-based measures of the staff's outlook in the spirit of \citet{aruoba2025} are the obvious candidate, and they
can be built from the same Bluebook and Greenbook corpus used here. Second, the $\lambda$-relaxation of effect
homogeneity can be paired with a relaxation of unconfoundedness along the lines of \citet{masten2018}. That would
report how much selection on unobservables is needed to overturn each sign, which is the question the results raise.
Third, the menus can be used with external shocks. \citet{kolesar2025} show that linear projections on an observed
shock identify a weighted average of causal effects even when the economy is nonlinear. Estimating such projections
within menus, with high-frequency surprises purged of prior news as in \citet{bauer2023nber}, would combine the
structural restriction documented here with a source of variation that does not rely on selection on observables.
